\section{Length multipliers and ray lifts}\label{sec:ray}

This section supplies the length estimate needed for the isolated
boundary factors. We construct ray lifts with uniform density and
intertwining bounds (Proposition~\ref{prop:ray-lift}), then combine them
with property~$(T)$ averaging to prove addition of rescaled spherical
length laws (Proposition~\ref{prop:ray-addition}). Addition will let us
compare repeated source segments in Section~\ref{sec:scale}, where it
rules out total superlinear growth and supplies the length window used
to construct the isolated heights.

The isolated components requiring the present estimate are real Lie
factors, as established in Sections~\ref{sec:transfer}
and~\ref{sec:finite}.  Their symmetric spaces are quaternionic
hyperbolic spaces of dimension at least two and the Cayley hyperbolic
plane.  Fix one such space \(X\).  Normalize its metric so that
\(r(g)=d(o,go)\) is the logarithm of the dilation in the indivisible
restricted root.  The argument uses an entire ambient factor; a
rank-one residue inside a larger factor does not supply this axis.

All other measured coordinates, including all finite places, remain
in the coefficient algebra.  Fourier indices remain distinct even
when their projections to \(\operatorname{Isom}(X)\) coincide.
The analytic estimates use only this real projection, whereas the
property-\((T)\) averaging below uses the whole ambient product or
its lattice.  No projected discreteness or one-factor counting
assumption enters these estimates.

Constants may depend on \(X\) and the specified smooth functions,
but not on the number of Fourier indices, scalar twists, or
coefficient-algebra dimensions.  The analytic estimates are stated
homogeneously in the operator and Hilbert norms, so arbitrary finite
coefficient traces are allowed there.  When an assertion is stated
uniformly along a sequence, every displayed Hilbert bound is part of
its hypotheses.  For a fixed finite algebra and fixed matrix size,
an operator bound implies the corresponding Hilbert bound.  All
varying finite factors used for spherical laws below have normalized
trace, and the quotient base has probability measure.

\subsection{Coordinates and a uniform smooth calculus}

Write the horospherical group as \(N_b=V\times Z\), with multiplication
\[
 (v,z)(v',z')=(v+v',z+z'+\tfrac12[v,v']),\qquad
 \delta_R(v,z)=(Rv,R^2z).
\]
The inner products can be chosen so that
\(\langle J_zv,w\rangle=\langle z,[v,w]\rangle\) and
\(J_z^2=-|z|^2\id\).  The horizontal layer is bracket-generating:
if \(z\) is orthogonal to all brackets, then \(J_z=0\), and the
displayed identity forces \(z=0\).  Set \(Q=\dim V+2\dim Z\).
In half-space coordinates \((n,\lambda)\in N_b\times(0,\infty)\), the
metric and volume, with a fixed normalization of Haar measure, are
\begin{equation}\label{eq:ray-metric}
 ds_X^2=\frac{d\lambda^2}{\lambda^2}
       +\frac{|dv|^2}{4\lambda^2}
       +\frac{|dz-\tfrac12[v,dv]|^2}{4\lambda^4},
 \qquad dp=\frac{dn\,d\lambda}{\lambda^{Q+1}}.
\end{equation}
These are the Iwasawa coordinates for the indicated symmetric spaces;
their two root spaces are precisely \(V\) and \(Z\). In the conventions
of \cite[Section~1, Proposition~1]{DamekRicci1992}, set $a=\lambda^2$
and divide the metric tensor by four to obtain
Equation~\eqref{eq:ray-metric}. In particular a vertical dilation has
length \(|\log R|\). The H-type description and the boundary geometry
are developed in \cite{CDKR1991,CDKR1998,Pansu1989}.

The horizontal term in \eqref{eq:ray-metric} determines the canonical
horizontal conformal class on \(\partial X\). The chart changes induced
by isometries preserve this class \cite[Lemma~9.6]{Pansu1989}; the metric
comparison in Lemma~\ref{lem:ray-geometry} below verifies this directly
in coordinates. Choose a smooth global
representative and average its pullbacks over the maximal compact
stabilizer of \(o\). All the averaged metrics lie in the same conformal
class, so their average is a metric in that class and is compact-invariant.
We use this metric throughout and let \(d_{\partial X}\) be its
Carnot--Carath\'eodory distance. In a buffered horospherical chart this
distance is comparable to \((|v|^4/16+|z|^2)^{1/4}\), where
\((v,z)=\log(n^{-1}n')\).

Here and below a scale box means a translate of the dilation by \(R\) of
a fixed bounded box in exponential coordinates, with a fixed larger box
available for smooth cutoffs.  On the compact boundary we use finitely
many such coordinate charts.  Bounded overlap and bounded adjacency refer
to constants for these fixed larger boxes.  A normalized \(C^m\) bound
is the ordinary \(C^m\) bound after translation and dilation to the fixed
box.  Auxiliary variables ranging in a fixed compact set may be appended.

\begin{lemma}[Distance and changes of scale]\label{lem:ray-geometry}
The following assertions hold with constants uniform on buffered compact
coordinate charts and, for every integer \(m\ge0\), in normalized
\(C^m\) norms.
\begin{enumerate}[label=(\roman*)]
\item If \((v,z)=\log(n^{-1}n')\), then
\begin{equation}\label{eq:ray-distance}
 \cosh d((n,\lambda),(n',\lambda'))
 =\frac{\sqrt{(\lambda^2+\lambda'^2+|v|^2/4)^2+|z|^2}}
 {2\lambda\lambda'}.
\end{equation}
\item Changes between the chosen compactly rotated horospherical
boundary charts preserve normalized bounds on boxes of every scale
\(0<R\le1\).  The same assertion holds in the
adjoined height variable \(0\le\lambda\le CR\).  The conversion between
the outward endpoint and \(e^{-r}\), and finite half-space coordinates,
has the form
\begin{equation}\label{eq:ray-collar}
 n=n_0\bigl(e^{-2r}U,e^{-4r}W\bigr),\qquad
 \lambda=e^{-r}b,
\end{equation}
where \(U,W,b\) are smooth in the endpoint and \(e^{-r}\), and \(b\)
is bounded above and below by positive constants.
\item Let \(g\) have length \(L\), with pole \(g^-\) and image endpoint
\(g^+\).  An input annulus of size \(q=e^{-u}\), where size includes
height, maps to an output annulus of size \(q'=e^{-L+u}\), for
\(0\le u\le L\).  On these annuli its normalized derivatives are
bounded.  A smaller boundary box of radius \(\delta\le cq\), whose
height is at most \(c\delta\), has image radius at most
\(C e^{2u-L}\delta\).  After normalization at input scale \(\delta\)
and output scale \(e^{2u-L}\delta\), all prescribed derivatives,
including the height derivatives, are uniformly bounded.  Exterior pieces
\(q\ge c\) have an output upper bound \(Ce^{-L}\).
\end{enumerate}
\end{lemma}

\begin{proof}
For completeness, the distance formula can be checked directly from the
metric.  Translate and dilate the first point to \((0,1)\), put
\(A=\lambda^2+1+|v|^2/4\), \(C=A^2+|z|^2\), and
\(c=\sqrt C/(2\lambda)\).  In the orthonormal frame of
\eqref{eq:ray-metric}, the three components of the gradient of
\(c\) are
\[
 \frac{\lambda A}{\sqrt C}-\frac{\sqrt C}{2\lambda},\qquad
 \frac{Av+J_zv}{2\sqrt C},\qquad
 \frac{\lambda z}{\sqrt C}.
\]
Since \(J_zv\perp v\) and \(|J_zv|^2=|z|^2|v|^2\), the sum of their
squared norms is \(c^2-1\).  Thus \(\operatorname{arcosh}c\) has gradient norm one
away from its unique zero \((0,1)\).  It is proper.  Its negative
gradient curves reach that zero in time \(\operatorname{arcosh}c\); conversely its
gradient bound bounds its variation by the length of any curve.  It is
therefore the distance from \((0,1)\), proving (i).

One convenient Weyl inversion, in these coordinates \cite[Section~2.4]{Pohl2010}, is
\begin{equation}\label{eq:ray-inversion}
 j(v,z,\lambda)=
 \left(\frac{(-A_0+J_z)v}{B_0},-\frac z{B_0},
             \frac\lambda{\sqrt{B_0}}\right),\quad
 A_0=\lambda^2+|v|^2/4,\quad B_0=A_0^2+|z|^2.
\end{equation}
The identity \(j\delta_R=\delta_{R^{-1}}j\) follows by substitution.
On a compact annulus defined by
\((|v|^4/16+|z|^2)^{1/4}+\lambda\asymp1\), the denominator is bounded
away from zero.  All derivatives of the displayed formula are therefore
bounded, including at height zero, and its height ratio is positive and
bounded.  These statements apply to the Iwasawa groups of the symmetric
spaces under consideration; the inversion is their Weyl action.  No
assertion about arbitrary nonsymmetric two-step groups is needed.

We spell out the coordinate issue in (ii).  In a finite Bruhat cell the
Gauss factors of a compact rotation times a horospherical translation
are smooth, with bounded derivatives on buffered compact sets.  Write
them as \(n'ma_b\bar n\).  Passing \(\bar n\) across \(a_{\log\lambda}\)
replaces its negative-root coordinates by \(\lambda Y_{-1}\) and
\(\lambda^2Y_{-2}\).  At \(o\) the first of these tangent directions
has no component in the central horospherical tangent space.  Taylor's
formula consequently gives horizontal coordinate \(O(\lambda)\) and
central coordinate \(O(\lambda^2)\), with smooth quotients, before the
outer dilation by \(b\lambda\).  After that dilation the two errors are
\(\lambda^2U\) and \(\lambda^4W\).  This proves
\eqref{eq:ray-collar}, including bounded derivatives and a
positive height ratio.

Boundary transitions are contact maps: a nonzero central image of a
horizontal tangent would, by \eqref{eq:ray-metric}, turn growth
of order \(\lambda^{-2}\) for the squared metric into growth of order
\(\lambda^{-4}\), contradicting invariance of the metric.  After
the contact conclusion, comparison of the leading
\(\lambda^{-2}\) terms shows that the horizontal derivative is a
positive dilation followed by an orthogonal map.  Thus the boundary
action is horizontal conformal.  After
recentring both charts, their central component therefore has no linear
horizontal term.  Substitute \((v,z,\lambda)=(RV,R^2Z,Rh)\) in Taylor's
formula.  The horizontal component is divisible by \(R\), the central
component by \(R^2\), and the height by \(R\), with bounded derivatives
of every prescribed order.  The same argument applied to the inverse
gives (ii).  The stronger errors in \eqref{eq:ray-collar} also
remain bounded after this substitution.

In compactly rotated pole and image charts, \(g\) is
\(\delta_{e^{-L}}j\), up to compact isotropy.  Rescale the input annulus
by \(q^{-1}\) and the output by \((q')^{-1}\).  The resulting map is
exactly the restriction of \(j\) to a fixed buffered annulus.
The derivative and height-ratio conclusions follow from
\eqref{eq:ray-inversion}.  Applying (ii) to a box of relative
radius \(\delta/q\) gives the last assertion of (iii).  A finite cover
away from the pole proves the exterior assertion.
\end{proof}

The horizontal metric also fixes the exact volume convention used
later. Let \(\mathcal V\) be the horizontal bundle and give
\(T\partial X/\mathcal V\) the quotient metric induced by the surjective
bracket map
\[
 \bigwedge^2\mathcal V\longrightarrow T\partial X/\mathcal V.
\]
The product of the horizontal and quotient volume densities is a smooth
density on \(\partial X\), independent of a splitting of its tangent
bundle. An isometry of \(X\) acts on \(\mathcal V_x\) by a dilation
\(d>0\) followed by an isometry. Naturality of the bracket makes its
induced map on the quotient a dilation \(d^2\) followed by an isometry.
Its Jacobian for this density is consequently
\[
 d^{\dim V}d^{2\dim Z}=d^Q.
\]
The density is compact-invariant. Since the maximal compact group is
transitive on the boundary, its normalization to mass one is exactly
the compact angular probability. Thus the horizontal dilation has
angular Jacobian equal to its \(Q\)th power. Here \(Q=\dim V+2\dim Z\)
is the homogeneous dimension; the ordinary boundary dimension is
\(\dim V+\dim Z\).

\begin{lemma}[Triangle defect]\label{lem:ray-defect}
There is a constant \(C\) such that, if \(x,y\in X\) and
\(b=d(o,x)+d(o,y)-d(x,y)>0\), their outward endpoints satisfy
\[
                   d_{\partial X}(x^+,y^+)\le Ce^{-b/2}.
\]
Consequently, if \(r(g)+r(h)-r(gh)\ge b>0\), then
\(d_{\partial X}(g^-,h^+)\le Ce^{-b/2}\).
\end{lemma}

\begin{proof}
Both radii are at least \(b/2\), by the triangle inequality.
Choose a compactly rotated finite chart containing both endpoints
in its buffered interior.  Finitely many such choices cover the
compact space of endpoint pairs, so their constants are uniform.
Let \(\delta=d_{\partial X}(x^+,y^+)\).  The heights of \(x,y\)
are comparable to their radial exponentials, and their boundary
coordinate errors are at most \(Ce^{-\min(d(o,x),d(o,y))}\), by
\eqref{eq:ray-collar}.  If \(\delta\) is at most a fixed
multiple of this quantity the assertion follows.  Otherwise the
homogeneous separation of their current-point horizontal coordinates
is comparable to \(\delta\), and
\eqref{eq:ray-distance} gives
\[
 d(x,y)\ge d(o,x)+d(o,y)+2\log\delta-C.
\]
Rearranging proves the assertion.  For the last statement apply it
to \(x=g^{-1}o\), \(y=ho\).
\end{proof}

\begin{lemma}[Smooth separation and cells]\label{lem:ray-cells}
Fix the dimensions of the row and column coordinate spaces and a bound
on the number of auxiliary compact variables.  For every prescribed
summable power \(N\), there are an integer \(m\) and a constant \(C\)
with the following properties.
\begin{enumerate}[label=(\roman*)]
\item A smooth symbol on the product of two fixed buffered coordinate
boxes, with \(C^m\) norm at most \(A\), has a separated Fourier
expansion with coefficients bounded by
\(CA(1+|\nu|+|\nu'|)^{-N}\).  Unary factors have absolute value at
most one and are supported in the respective larger boxes.  This also
holds for unary symbols, with an arbitrary external parameter in their
coefficients.
\item If the pairs of boxes supporting a kernel form a graph of degree
at most \(b\), and the normalized extensions in (i) have bounds \(A\),
the kernel is a Schur multiplier of norm at most \(C_b A\).
Its matrix amplifications have the same bound.  The number of boxes is
irrelevant.
\end{enumerate}
\end{lemma}

\begin{proof}
Multiply an extension by fixed cutoffs and view it as a smooth function
on a fixed torus.  Integration by parts with
\((1-\Delta)^k\) gives (i) whenever \(2k\ge N\); choose \(N\) larger
than the dimension plus two.  At a later use requiring derivatives of
the unary factors, increase \(N\) by the required derivative order.
Thus every assertion here uses a finite, specified number of
derivatives, chosen before any scale tends to infinity.

For a fixed Fourier term and a matching of row and column boxes, the
kernel is a sum of disjoint rectangular rank-one kernels and has Schur
norm bounded by the coefficient bound times the overlap constant.
A graph of degree \(b\) is the union of at most \(2b-1\) matchings:
greedily color its edges, first for finite subgraphs and then pass to a
subsequence of the finite colorings.  Split boundedly overlapping
supports by assigning each point to one of its boundedly many boxes;
the assignment is a row or column multiplier of norm one.  Summing the
matchings and the absolutely convergent Fourier series proves (ii).
Each rank-one estimate is a factorization estimate and is unchanged
under matrix amplification.
\end{proof}

\subsection{Uniform radial multipliers}

Completely bounded radial approximations on these rank-one groups are
classical \cite{CowlingHaagerup1989}. The fixed profiles and uniform
coefficient-algebra estimates used below require the explicit smooth
separation argument that follows.
The search for a uniform fixed-profile radial Schur estimate was
motivated in part by Ozawa's bounded-degree hyperbolic-graph estimates
\cite[Theorem~1]{Ozawa2008}. Those graph-metric estimates are not
applied here; the exact continuous kernel is treated by the Carnot-cell
argument below.

\begin{proposition}[Radial cutoffs]\label{prop:ray-radial}
Let \(f\in C^\infty([0,\infty))\) be constant on a neighborhood of zero
and on a neighborhood of infinity.  There are \(R_f,C_f>0\) such that
the kernel
\[
 (x,x')\longmapsto f(d(x,x')/R)
\]
has Schur multiplier norm at most \(C_f\) for every \(R\ge R_f\).
Consequently, if a countable group \(\Gamma\) has any homomorphism into
\(\operatorname{Isom}(X)\), \(\eta\) is a normalized scalar cocycle,
and \(B\) is a finite tracial von Neumann algebra, then
\[
 D_{f(r/R)}:\ B\,\bar\otimes L_\eta(\Gamma)\longrightarrow
 B\,\bar\otimes L_\eta(\Gamma),\qquad
 \sum_hY_hu_h\longmapsto\sum_h f(r(h)/R)Y_hu_h
\]
has completely bounded norm at most \(C_f\).  The constants are
independent of \(\Gamma,\eta,B\) and of the homomorphism.
\end{proposition}

\begin{proof}
Subtract the constant value of \(f\) near infinity.  Choose \(0<a<b\)
such that \(f\) is constant on \([0,2a]\) and vanishes on \([b,\infty)\).
Write \(s=-\log\lambda\), \(s'=-\log\lambda'\).  Since
\(|s-s'|\le d(x,x')\), only a fixed number of offset bands between
height blocks of length \(R\) can contribute.  At a fixed offset the
block pairs are a matching.  Simultaneous dilation permits us to work
with \(|s|,|s'|\le C_bR\) in each pair.

Take smooth radial cutoffs \(\chi_j(n^{-1}n')\) at radii
\(\epsilon_j=e^{-j}\), equal to one on a smaller ball, and choose
\(J=\lceil C'_bR\rceil\), with \(C'_b>C_b+b/2+2\).  Formula~\eqref{eq:ray-distance}
then implies that \(\chi_{-J}=1\) throughout the support of
\(f(d/R)\).  The annuli
\(\chi_j-\chi_{j+1}\), \(-J\le j<J\), and the innermost ball
\(\chi_J\) cover the support of the kernel.  On each annulus set
\[
 H_j=\log\frac{\max(\epsilon_j,\lambda,\lambda')^2}
                         {\lambda\lambda'}.
\]
The following estimates are uniform in \(j,R\): the kernel
\(f(d/R)-f(H_j/R)\), restricted to that annulus, has Schur norm
\(O(R^{-1})\); the same holds on the innermost ball; and the height
kernel \(f(H_j/R)-f(H_{j-1}/R)\) has Schur norm \(O(R^{-1})\).
We verify the details, since an uncontrolled height transition would
invalidate the assertion.

Put \((v,z)=\delta_{\epsilon_j}(V,Z)\).  On an annulus \((V,Z)\)
lies in a fixed compact set separated from zero.  Each ratio involving
only one height and the fixed scale \(\epsilon_j\) is adjoined as an
independent compact unary coordinate before Fourier separation;
derivatives in \(s/R,s'/R\) keep these auxiliary coordinates fixed.  If
\(\lambda,\lambda'\le\epsilon_j\), put
\(\alpha=\lambda/\epsilon_j\), \(\beta=\lambda'/\epsilon_j\).
Equation~\eqref{eq:ray-distance} gives
\[
 2\cosh d=e^{H_j}\sqrt P,\qquad
 P=(\alpha^2+\beta^2+|V|^2/4)^2+|Z|^2.
\]
Here \(P\) is bounded above and below by positive constants.  If
\(\lambda\ge\epsilon_j\ge\lambda'\), instead put
\(\alpha=\epsilon_j/\lambda\), \(\beta=\lambda'/\epsilon_j\);
the same formula holds with
\[
 H_j=\log(\lambda/\lambda'),\qquad
 P=(1+(\alpha\beta)^2+\alpha^2|V|^2/4)^2+\alpha^4|Z|^2.
\]
Again \(P\) and all its derivatives are bounded and \(P\ge1\).
The reversed mixed case is identical.  The four cuts made so far are
rectangular height cuts; their Schur norms are one.

If both heights exceed \(\epsilon_j\), distance differs by a bounded
amount from \(|s-s'|\).  Where \(|s-s'|\le aR\), both arguments of
\(f\) are in its constant interval for sufficiently large \(R\).
Outside this interval use smooth sign cutoffs in \((s-s')/R\).
On the positive piece, for example, take the larger height as the
normalizing scale.  The formula for \(P\) is the preceding one with
the ratio of the smaller height to the larger height in place of
\(\alpha\beta\).  In the bounded variables \(s/R,s'/R\) that ratio
is \(e^{-R|s/R-s'/R|}\), separated from its nonsmooth transition at
zero.  Every fixed derivative is uniformly bounded, because
\(R^m e^{-aR}\) is bounded.  The sign cutoffs themselves have uniformly
bounded derivatives and hence bounded Schur norm by smooth separation.

In each of these cases, when \(H_j\ge aR\), write explicitly
\[
 d-H_j=\tfrac12\log P+
 \log\left(\frac{1+\sqrt{1-4e^{-2H_j}/P}}2\right).
\]
The last function is evaluated in a compact interval strictly inside
its domain of smoothness.  It has bounded derivatives in the compact
ratio variables just described and in \(s/R,s'/R\).  Where \(H_j\)
is smaller the multiplier difference is zero, after a smooth
transition still inside the constant interval of \(f\).  Thus
\[
 f(d/R)-f(H_j/R)
 =\frac{d-H_j}{R}\int_0^1
             f'\bigl(H_j/R+t(d-H_j)/R\bigr)\,dt
\]
has normalized \(C^m\) norm at most \(C_m/R\).  Use adjacent
\(\epsilon_j\)-cells for \(n,n'\).  Their relative coordinates
\((V,Z)\) have uniform bounds.  Lemma~\ref{lem:ray-cells}
gives the required Schur estimate.  The innermost ball requires only
the large-height cases: \(J\) was chosen larger than all the height
exponents.  Past the exterior cutoff the distance is at least \(bR\).

For the height increments use
\[
 H_j=s+s'-2\min(j,s,s').
\]
Split each of \(s,s'\) into \(( -\infty,j-1]\), \([j-1,j]\),
and \([j,\infty)\).  If either is in the first interval, the
increment vanishes.  If both are in the last interval, \(H_j\)
and \(H_{j-1}\) differ by exactly two; their scaled \(f\)-difference
has \(C^m\) norm \(O(R^{-1})\) in \(s/R,s'/R\).  If exactly one
variable lies in the middle interval, use its unit-scale coordinate
and the other variable divided by \(R\), with the same conclusion.
If both lie in the middle interval, both arguments of \(f\) are in
its constant interval and the difference is zero.  These are
rectangles, so no triangular projection on height indices occurs.

Finally write \(\phi_j=f(H_j/R)\).  The exact telescoping identity is
\[
 \sum_{j=-J}^{J-1}\phi_j(\chi_j-\chi_{j+1})+\phi_J\chi_J
 =\phi_{-J}\chi_{-J}
       +\sum_{j=-J+1}^J(\phi_j-\phi_{j-1})\chi_j.
\]
Our choice of \(C'_b\) also gives \(H_{-J}>bR\), so
\(\phi_{-J}=0\).  Each ball kernel has uniformly bounded Schur norm:
its two horizontal entries belong to adjacent scale cells and its
normalized symbol is smooth.  There are \(O(R)\) increments and
\(O(R)\) annular errors, each costing \(O(R^{-1})\).  Summing the
fixed number of height-block bands proves the Schur bound.

Restrict the kernel to the orbit points \(h^{-1}o\).  The regular
matrix of the Fourier multiplier is its entrywise multiplication
kernel, up to the scalar phases already present in the twisted
regular representation.  The factorization proof of the Schur bound
is unchanged by these phases or by tensoring with \(B\).
Coincident orbit points cause no difficulty.  This proves the last
assertion, including complete boundedness.
\end{proof}

\subsection{One-sided localization}

Let \(D=B\bar\otimes L_\eta(\Gamma)\), with its product trace \(\tau\).
Write \(Y_h\in L^2(B)\) for the Fourier coefficients of \(Y\).  For a
bounded scalar function \(\varphi\) on \(\Gamma\), \(D_\varphi\) denotes
coefficient multiplication on \(L^2(D)\); it need not be bounded on
\(D\).  Phases in the coefficient formula for \(YY^*\) have absolute
value one.  For finite Fourier sums that formula is
\[
 (YY^*)_g=\sum_h
 \eta(gh,h^{-1})\overline{\eta(h,h^{-1})}\,Y_{gh}Y_h^*.
\]
All coefficient identities below are first understood for such sums.

\begin{lemma}[Unary Plancherel estimate]\label{lem:ray-unary}
Suppose \(\varphi_{i,\nu}:\Gamma\to\mathbb C\), with \(i\) in a
countable set, satisfy
\[
 \sup_h\sum_i|\varphi_{i,\nu}(h)|^2\le b
\]
for every \(\nu\).  Let \(T_g\) be an arbitrary family of indices, and
associate \(i(g,t)\) to each \(t\in T_g\), with at most \(b\) occurrences
of any given \(i\) for a fixed \(g\).  Suppose
\[
 z_{g,t}=\sum_\nu c_{g,t,\nu}
       [Y(D_{\varphi_{i(g,t),\nu}}Y)^*]_g,
 \qquad |c_{g,t,\nu}|\le A a_\nu,
 \qquad \sum_\nu a_\nu<\infty.
\]
Then
\begin{equation}\label{eq:ray-unary}
 \left(\sum_g\sum_{t\in T_g}\|z_{g,t}\|_{2,B}^2\right)^{1/2}
 \le b A\left(\sum_\nu a_\nu\right)
               \|Y\|_\infty\|Y\|_2.
\end{equation}
There is an identical estimate when the localized factor is the first
factor of \(YY^*\).  In particular coefficients, admissibility of a
cell, and the choice of finitely many companion cells may depend on
\(g\) without affecting the estimate.
\end{lemma}

\begin{proof}
For one \(\nu\), discard the admissibility restriction after using
the multiplicity bound.  Plancherel and the operator norm bound for
the uncut factor give
\[
 \begin{split}
 \sum_{g,t}\|[Y(D_{\varphi_{i(g,t),\nu}}Y)^*]_g\|_2^2
 &\le b\sum_i\|Y(D_{\varphi_{i,\nu}}Y)^*\|_2^2\\
 &\le b\|Y\|_\infty^2
                   \sum_i\|D_{\varphi_{i,\nu}}Y\|_2^2
 \le b^2\|Y\|_\infty^2\|Y\|_2^2.
 \end{split}
\]
Minkowski's inequality sums the absolutely summable frequencies.  Taking
adjoints proves the other orientation.  Notice that no operator norm
for a localized column was used.
\end{proof}

\begin{lemma}[Angular column estimate]\label{lem:ray-angular}
There are constants \(c,C>0\) and an integer \(m\), depending only on
\(X\), with the following property.  Put \(\delta=e^{-s}\), \(s\ge1\).
Let \((f_C)_C\) be smooth masks supported on buffered boundary boxes of
radius \(\delta\), with fixed overlap and normalized \(C^m\) bounds.
If \(Y\in D\) has Fourier support in \(\{r\ge l\}\), where
\(l\ge Cs+C\), then, with \(A_C=D_{f_C(h^+)}Y\),
\begin{equation}\label{eq:input-7}
 \sum_{C,D'}\|A_CA_{D'}^*\|_2^2
 =\left\|\sum_C A_C^*A_C\right\|_2^2
 \le C\|Y\|_\infty^2\|Y\|_2^2.
\end{equation}
The constants can depend on the stated mask bounds, but not on their
number.  The assertion holds with arbitrary finite tracial coefficient
algebras and with any scalar twist.
\end{lemma}

\begin{proof}
We estimate the coefficients of \(A_CA_{D'}^*\) at a fixed shift \(g\).
Put \(L=r(g)\).  Their weights in the coefficient formula for \(YY^*\)
are
\begin{equation}\label{eq:ray-pair-mask}
             f_C((gh)^+)\overline{f_{D'}(h^+)}.
\end{equation}
Only pairs with both \(h\) and \(gh\) in the long support matter.
Their polar heights are at most \(e^{-l}\).  We describe the smooth
extensions of these weights; all coordinate boxes below come from
fixed covers independent of \(g\).

First suppose that the input box is separated from \(g^-\) by at
least a sufficiently large multiple of \(\delta\), and its pole
distance is \(q=e^{-u}\).  If \(2u\le L+C_1\),
Lemma~\ref{lem:ray-geometry}(iii) says that the map from this
box to the output boundary, including the extra coordinate
\(e^{-r(h)}/\delta\), has normalized derivatives bounded by a constant.
Its image meets only boundedly many output boxes.  Extend on a fixed
larger box, multiply by its buffered cutoff, and Fourier-expand there.
The weight in \eqref{eq:ray-pair-mask} is now a unary expansion
in \(h\), with coefficients decaying faster than any prescribed fixed
power.  The basis functions depend on \(D'\), the fixed coordinate
chart, and the frequency, and not on \(g\) or \(C\).  If
\(2u>L+C_1\), inversion gives output pole distance
\(e^{-L+u}\); unless the output box is itself a pole box, the inverse
map has the same bounded-magnification property.  Use the reversed
unary estimate in that case.  A box intersecting the transition
\(2u=L+O(1)\) admits either construction, with a larger fixed buffer.
Starting from an output box gives the same alternatives.  Bounded
\(L\) is included by increasing the constants.

Thus every pair except pairs of boxes lying within \(C_2\delta\)
of the two respective poles is handled by
Lemma~\ref{lem:ray-unary}.  To justify the multiplicity used
there, a map with bounded magnification sends one buffered
\(\delta\)-box into a ball of radius \(C\delta\), which meets a
bounded number of reference boxes.  The height variable changes this
radius by at most \(Ce^{-l}\).  This proves the required multiplicity
for each \(g\) and initial box, without a restriction on the number of
group elements in it.

There remain only boundedly many exceptional pairs for each \(g\).
If such a pair is compatible with the long supports, the two
current-point pole sizes have product comparable to \(e^{-L}\), by
\eqref{eq:ray-inversion}, whereas each is at most
\(C\delta\).  Consequently
\begin{equation}\label{eq:ray-long-shift}
                         L\ge2s-C_3.
\end{equation}
Expand each mask as its value at its pole plus its difference from
that value.  The constant term costs at most
\(C\|YY^*\|_2\).  Terms with just one nonconstant mask are covered
by Lemma~\ref{lem:ray-unary}; any additional constant term
created by this expansion is treated with the first term.

For the product of the two differences, partition the input
current-point pole coordinate, including height, into smooth
unit-width annuli with size \(e^{-u}\),
\(0\le u\le L/2+C_4\).  Use the analogous output annuli on the
remaining part.  A smooth transition at the scale \(e^{-L/2}\)
has bounded normalized derivatives on both descriptions, by
\eqref{eq:ray-inversion}.  Exterior annuli mean a fixed
finite cover away from the pole.

On an input annulus, the two differences, after substitution, have
product seminorm bounded by
\begin{equation}\label{eq:ray-annular-decay}
 C_m\min(1,e^{s-u})\min(1,e^{s-L+u})
                         \le C'_m e^{-|s-u|}.
\end{equation}
Here is the derivative justification for both factors.  On scales
smaller than a mask box, its difference from its pole value is
bounded by the relative scale, by Taylor's formula; the same bound
holds for every normalized derivative on the smaller box.  On scales
larger than the mask box by a fixed factor, the mask is zero on the
long support, because its current-point height is much smaller than
\(\delta\).  It can therefore be replaced by zero on that part before
subtracting its pole value, giving a bounded constant symbol.  At
comparable scales its given seminorm bound suffices.  The output
factor has exactly the same three descriptions after the normalized
inversion, with scale \(e^{-L+u}\).  The collar substitutions have
bounded derivatives including the normalized height by
Lemma~\ref{lem:ray-geometry}(ii).  These constructions give
smooth extensions on buffered boxes; equality is required only on
the two long supports.  Finally, if \(u\le s\), use
\eqref{eq:ray-long-shift} in the second minimum; if \(u\ge s\),
use the first minimum.  This proves the last inequality in
\eqref{eq:ray-annular-decay}.

At a fixed integer depth \(u\), use the reference boundary boxes of
radius \(e^{-u}\), and append the height divided by \(e^{-u}\) as
an auxiliary unary variable.  Only boundedly many such boxes meet
the annulus about the given pole.  Their basis functions are chosen
before \(g\); the coordinates of the pole merely change the
expansion coefficients.  The number of original exceptional pairs
is also bounded.  Lemmas~\ref{lem:ray-cells} and
\ref{lem:ray-unary} therefore give a bound
\(Ce^{-|s-u|}\|Y\|_\infty\|Y\|_2\) in the square sum of all
coefficients using this depth.  Sum over \(u\) in norm; the geometric
series has a uniform sum.  The reversed annuli give the same bound.
Combining the nonexceptional and exceptional estimates proves the
inequality in \eqref{eq:input-7}.  Its equality is the trace
identity
\(\tau(A_CA_{D'}^*A_{D'}A_C^*)
=\tau(A_C^*A_CA_{D'}^*A_{D'})\), summed over the masks.

For general \(Y\), choose bounded Fourier-polynomial approximations
in the strong-star topology and apply a smooth radial cutoff equal
to one on \(\{r\ge l\}\) and zero on \(\{r\le l/2\}\).
Proposition~\ref{prop:ray-radial} bounds these approximants
uniformly.  Increase the constant in the hypothesis on \(l\) to
allow this slack.  Each masked column converges in \(L^2\), so each
product of two masked columns converges in \(L^1\), by tracial
Cauchy--Schwarz.  For a fixed finite family of pairs, the preceding
estimate bounds their products in the Hilbert direct sum of
\(L^2(D)\).  Weak compactness gives a weak \(L^2\) limit, and the
\(L^1\) convergence identifies it with the desired products.
Lower semicontinuity proves the bound for this finite family;
exhaustion proves it for all pairs.  In particular the diagonal
terms put every masked column in \(L^4\), so the trace identity in
\eqref{eq:input-7} remains valid.  This also explains why neither infinite
Fourier support nor a nondiscrete projected group changes the proof.
\end{proof}

\subsection{The ray lift and its densities}

Let \(\mathcal K_X=L^2(X,dp)\), and let \(T_g\) be translation on
\(\mathcal K_X\).  Define
\[
 \Delta:D\longrightarrow B(\mathcal K_X)\bar\otimes D,
 \qquad \Delta(bu_h)=T_h\otimes bu_h.
\]
This is a normal unital homomorphism.  Indeed the unitary on
\(\mathcal K_X\otimes\ell^2\Gamma\) that sends
\(\zeta\otimes\delta_h\) to \(T_h^{-1}\zeta\otimes\delta_h\)
conjugates \(T_g\otimes u_g\) to \(1\otimes u_g\), with the same
scalar twist.  Tensoring with the coefficient representation proves
normality in the displayed algebra.
The right action of \(D\) on \(\mathcal K_X\otimes L^2(D)\) acts on its
second factor.  Choose real smooth functions \(\psi\) and \(F\), with
\(\supp\psi\Subset(1,2)\), \(F\) compactly supported on \(N_b\) and
invariant under the compact isotropy of a ray, such that
\(\int\psi^2=\int F^2=1\).  For \({d_1}\ge1\), a basepoint \(x\), and an
ideal endpoint \(v\), use a chart sending \(v\) to infinity and put
\begin{equation}\label{eq:ray-function}
 \theta_{x,d_1}(p,v)={d_1}^{-1/2}
 \psi\left(\frac{\log(\lambda_p/\lambda_x)}{d_1}\right)
 F\bigl(\delta_{\lambda_p^{-1}}(n_x^{-1}n_p)\bigr).
\end{equation}
Changing that chart by a horospherical translation, dilation, or
compact ray isotropy leaves the formula unchanged.  It is therefore
equivariant in \(x,p,v\).  The substitution
\(n_p=n_x\delta_{\lambda_p}q\), \(t=\log(\lambda_p/\lambda_x)/{d_1}\)
shows that \(\int\theta_{x,d_1}(p,v)^2dp=1\).
Its support stays a bounded distance from the ray from \(x\) to
\(v\), at depths between \({d_1}-O(1)\) and \(2{d_1}+O(1)\).
The coefficient isometry is
\[
 \mathbb V_{d_1}Y=\sum_h\theta_{o,d_1}(\,\cdot\,,h^+)\otimes Y_hu_h.
\]
At a zero-displacement index choose any fixed direction; none of the
long-support assertions below depends on that choice.

For a column \(\xi\in\mathcal K_X\otimes L^2(D)\), its left and
right densities are the positive \(L^1(D)\) elements representing
the vector functionals for the actions \(\Delta(D)\) and
\(D^{\mathrm{op}}\), respectively.  Thus
\[
 \langle\xi,\Delta(x)\xi\rangle=\tau(\rho^L_\xi x),\qquad
 \langle\xi,\xi x\rangle=\tau(\rho^R_\xi x),\qquad
 \rho^R_\xi=\xi^*\xi.
\]
The density on the whole amplified left algebra is \(\xi\xi^*\),
with its product semifinite trace.  The following proposition puts
all three densities in trace \(L^2\).

\begin{proposition}[Ray lift]\label{prop:ray-lift}
Consider any sequence of finite tracial coefficient algebras, twists,
and countable group homomorphisms into \(\operatorname{Isom}(X)\).
Let \({d_1}\to\infty\), \(l/{d_1}\to\infty\), and let \(Y\in D\) have Fourier
support in \(\{r\ge l\}\).  Then the following statements hold.
For \textup{(ii)}, assume that
\(\|x\|_\infty,\|x\|_2,\|Y\|_\infty,\|Y\|_2\) remain bounded;
its convergence is uniform under these bounds.
For \textup{(iii)}, assume only the three norm bounds stated there.
\begin{enumerate}[label=(\roman*)]
\item For all sufficiently large indices, the vector \(\mathbb V_{d_1}Y\)
has left density on \(D\), for the
action \(\Delta\), and right density on \(D^{\mathrm{op}}\), both in
\(L^2\).  Their norms are at most
\(C\|Y\|_\infty\|Y\|_2\).  Also
\(\|\mathbb V_{d_1}Y(\mathbb V_{d_1}Y)^*\|_{2,\Tr\otimes\tau}\)
has this bound.
\item If \(\supp\widehat x\subset\{r\le a\}\), \(a/{d_1}\to0\), then
\[
 \|\mathbb V_{d_1}(xY)-\Delta(x)\mathbb V_{d_1}Y\|_2\longrightarrow0,
 \qquad
 \|\mathbb V_{d_1}(Yx)-(\mathbb V_{d_1}Y)x\|_2\longrightarrow0.
\]
The lift in these formulas uses the same \({d_1}\); the products have
support in \(\{r\ge l-a\}\).
\item Suppose \(Z\in L^2(D)\) is another column with Fourier support
in \(\{r\geq l\}\), and \(x\in D\) is bounded.
Suppose that for every contributing triple
\(h\in\supp\widehat Y\), \(hk\in\supp\widehat Z\),
\(k\in\supp\widehat x\), one has
\[
 d_{\partial X}(h^+,(hk)^+)\le e^{-c l'},
 \qquad c>0,\qquad l'/{d_1}\longrightarrow\infty.
\]
Then
\begin{equation}\label{eq:ray-close-pairing}
 \langle\mathbb V_{d_1}Z,(\mathbb V_{d_1}Y)x\rangle
                         -\langle Z,Yx\rangle\longrightarrow0.
\end{equation}
The assertion holds uniformly with
\(\|x\|_\infty,\|Y\|_2,\|Z\|_2\) bounded.
\end{enumerate}
All assertions remain valid for finite matrix versions and for
measurable fields, with the indicated operator bounds uniform and the
indicated Hilbert bounds taken in the integrated tracial norms.
\end{proposition}

We first prove the right-density bound in (i), then establish the overlap
estimate and the left-density bound, and finally prove (iii) followed by
(ii).

\begin{proof}[Proof of the right-density assertion]
Write
\(K_{d_1}(v,w)=\int\theta_{o,d_1}(p,v)\theta_{o,d_1}(p,w)\,dp\).
Disintegrate this integral at radial depth \(s=d(o,p)\).  Its range
is \([{d_1}-C,2{d_1}+C]\).  With measure \(ds/{d_1}\), the resulting kernels
\(k_{d_1,s}(v,w)\) are supported on adjacent boundary boxes of radius
\(e^{-s}\), with normalized \(C^m\) bounds independent of \({d_1},s\),
for every fixed \(m\).

Here are the normalization details.  The polar volume is
\(J(s)\,ds\,d\nu\), where
\(J(s)=c(\sinh s)^{\dim V}(\sinh2s)^{\dim Z}\asymp e^{Qs}\)
for \(s\ge1\).  The shadow of one ray has angular volume
\(O(e^{-Qs})\).  Formula~\eqref{eq:ray-function} contributes
\({d_1}^{-1}\), removed by the choice of \(ds/{d_1}\).  Thus no volume
factor remains.  To check derivatives, choose smooth compact frames
near one endpoint, translate back along its ray by depth \(s\), and
rescale neighboring directions by \(e^{-s}\).  Their negative-root
coordinates have sizes \(e^{-s}\) and \(e^{-2s}\); conjugation by
the split translation makes both bounded.  The positive-root
coordinates contract and the remaining factors stay compact.
Lemma~\ref{lem:ray-geometry}(ii), followed by
\eqref{eq:ray-function} on this fixed bounded set, therefore
bounds every prescribed normalized derivative.  The same argument
applies to the polar Jacobian after its factor \(e^{Qs}\) is removed.

Separate \(k_{d_1,s}\) on adjacent boxes.  At each frequency the right
density is a sum \(\sum_{C,D'}t_{C,D'}A_C^*B_{D'}\), where
\((t_{C,D'})\) is a scalar matrix of bounded degree and bounded
operator norm, and both mask families satisfy the hypotheses of
Lemma~\ref{lem:ray-angular}.  The column form of tracial
H\"older gives
\[
 \left\|\sum_{C,D'}t_{C,D'}A_C^*B_{D'}\right\|_2
 \le \|t\|
 \left\|\sum_C A_C^*A_C\right\|_2^{1/2}
 \left\|\sum_{D'}B_{D'}^*B_{D'}\right\|_2^{1/2}.
\]
Equation~\eqref{eq:input-7} bounds this by
\(C\|Y\|_\infty\|Y\|_2\).  Fourier coefficients are summable;
choose their decay order larger than the polynomial derivative
growth of the masks.  Integration in \(ds/{d_1}\), over an interval of
bounded total mass, proves the right-density bound.  Finally
\[
 \|\mathbb V_{d_1}Y(\mathbb V_{d_1}Y)^*\|_2^2
 =\| (\mathbb V_{d_1}Y)^*\mathbb V_{d_1}Y\|_2^2
\]
by the semifinite trace, first for finite-rank columns and then by
monotone approximation.  This proves the additional assertion in (i).
\end{proof}

\begin{lemma}[Common-endpoint overlap]\label{lem:ray-overlap}
For a displacement \(g\) of length \(L\), put
\[
 w_{d_1}(g,v)=\int_X\theta_{o,d_1}(p,v)\theta_{go,d_1}(p,v)\,dp.
\]
There are smooth annular partitions about \(g^+\), at integer depths
\(0\le j\le L\), with an exterior piece and an innermost ball,
such that on each piece
\begin{equation}\label{eq:ray-overlap-step}
 \begin{aligned}
 \|w_{d_1}(g,\cdot)-a_{d_1}(j,L)\|_{C^m_{\rm normalized}}
 &\le C_m/{d_1},\\
 a_{d_1}(j,L)&=\frac1{d_1}\int_{s>j}\psi(s/{d_1})
                  \psi((s-2j+L)/{d_1})\,ds.
 \end{aligned}
\end{equation}
The constants are uniform in \(L,{d_1}\ge1\).  Further,
\begin{equation}\label{eq:ray-overlap-increments}
 |\partial_j a_{d_1}(j,L)|\le C/{d_1},\qquad
 a_{d_1}(0,L)=1+O\bigl(\min(1,L/{d_1})\bigr).
\end{equation}
Noninteger endpoints of the last annulus are absorbed in its fixed
buffer.  When \(L\) is bounded, a fixed finite cover suffices.
\end{lemma}

\begin{proof}
We first prove uniformity in all angular derivatives, including
at the two ends of the annular decomposition.  By a compact rotation
put \(go=(0,\epsilon)\), \(\epsilon=e^{-L}\), and \(g^+=0\).
For a finite endpoint \(v=(V,W)\), define
\[
 p(v)=|V|^2/4,\quad D_h(v)=(h^2+p(v))^2+|W|^2,\quad
 P_h(v)=\left(\frac{(h^2+p(v)+J_W)V}{D_h(v)},
                         \frac W{D_h(v)}\right).
\]
Formula~\eqref{eq:ray-inversion} shows that
\(\delta_{\sqrt{D_1(v)}}L_{P_1(v)^{-1}}jL_{v^{-1}}\) fixes \(o\)
and sends \(v\) to infinity.  In this chart the displaced point is
\((n,e^\beta)\), with
\begin{equation}\label{eq:ray-endpoint-coordinates}
 n=\delta_{\sqrt{D_1(v)}}(P_1(v)^{-1}P_\epsilon(v)),
 \qquad e^\beta=\epsilon\sqrt{D_1(v)/D_\epsilon(v)}.
\end{equation}
All products in this formula are products in the real two-step
group \(N_b\); no nonassociative scalar multiplication is used.

Put \(v=\delta_Ru\), \(R=e^{-j}\), \(t=\epsilon/R\), on a fixed
buffered annulus for \(u=(V,W)\), and write
\[
 \begin{split}
 \Delta_R(u)&=(1+R^2p(u))^2+R^4|W|^2,\\
 Q_R(u)&=\left(
 \frac{R^2(1+R^2p(u)+R^2J_W)V}{\Delta_R(u)},
 \frac{R^4W}{\Delta_R(u)}\right).
 \end{split}
\]
Direct substitution in \eqref{eq:ray-endpoint-coordinates}
gives
\begin{equation}\label{eq:ray-endpoint-scaled}
 \beta=2j-L+B(R,t,u),\quad
 B=\tfrac12\log(\Delta_R/D_t),\quad
 \delta_Rn=m(R,t,u)
   =\delta_{\sqrt{\Delta_R}}(Q_R^{-1}P_t).
\end{equation}
For \(R,t\in[0,1]\), the denominator \(D_t\) is bounded away from
zero on the buffered annulus, and \(\Delta_R\ge1\).  Thus \(B,m\)
have bounded derivatives of every fixed order.  As \(R,t\to0\),
\(m\to P_0(u)\), uniformly with derivatives; this limit stays away
from the identity.  Consequently \(m\) is uniformly separated
from the identity if both \(j\) and \(L-j\) exceed a fixed constant.
For the innermost ball take \(j=L,t=1\): \(D_1\ge1\), so all upper
bounds remain valid, including at \(u=0\).  The exterior follows
from finitely many compact charts away from the pole, with
\(\beta=-L+O_{C^m}(1)\), \(n=O_{C^m}(1)\).

Let \(H(m)=\int F(q)F(m^{-1}q)\,dq\).  Then \(H\) is smooth and
compactly supported, and \(H(e)=1\).  Haar substitution in
\eqref{eq:ray-function} gives the exact formula
\begin{equation}\label{eq:ray-overlap-exact}
 w_{d_1}(g,v)=\int\psi(t)\psi(t-\beta/{d_1})
                         H(\delta_{e^{-t d_1}}n)\,dt.
\end{equation}
On a normalized annulus it becomes
\[
 \frac1{d_1}\int\psi(s/{d_1})\psi((s-2j+L-B(u))/{d_1})
                         H(\delta_{e^{j-s}}m(u))\,ds.
\]
Replace \(B\) by zero using the mean-value formula.  This costs
\(O_m({d_1}^{-1})\) after angular differentiation: the extra factor
is \({d_1}^{-1}\), and all derivatives of the accompanying \(H\)
factor are bounded on the relevant support.  Indeed, in the interior
annuli, compact support of \(H\) and separation of \(m\) from the
identity force \(s-j\ge-C\); near \(j=0\), the first \(\psi\)
forces \(s-j\ge {d_1}-C\); near \(j=L\), either of the shifted
\(\psi\)'s in the mean-value formula forces the same inequality.
The interval of integration has length \(O({d_1})\).

It remains to replace \(H(\delta_{e^{j-s}}m(u))\) by
\(\mathbf1_{s>j}\).  On the interior annuli, with \(r=s-j\),
every angular derivative satisfies the integrable bound
\[
 \left|\partial_u^\gamma
 \left(H(\delta_{e^{-r}}m(u))-\mathbf1_{r>0}\right)\right|
 \le C_\gamma\left(
       \mathbf1_{[-C,0]}(r)+e^{-r}\mathbf1_{[0,\infty)}(r)\right).
\]
For negative \(r\) this follows from the support of \(H\); for
positive \(r\) it follows from Taylor's formula at the identity,
because horizontal and central coordinates and all their
\(u\)-derivatives acquire factors \(e^{-r}\) and \(e^{-2r}\).
Integration against \(ds/{d_1}\) gives \(O_m({d_1}^{-1})\).  In the two
end bands the \(\psi\) supports already imply \(r\ge {d_1}-C\), so
the same conclusion follows without a lower bound on \(m\).
This proves \eqref{eq:ray-overlap-step} on all pieces.

Finally differentiation under the integral gives
\[
 \partial_ja_{d_1}(j,L)=-\frac1{d_1}\psi(j/{d_1})\psi((L-j)/{d_1})
 -\frac2{{d_1}^2}\int_j^\infty\psi(s/{d_1})
                              \psi'((s-2j+L)/{d_1})\,ds.
\]
This is \(O({d_1}^{-1})\).  At \(j=0\) the integral is
\(\int\psi(t)\psi(t+L/{d_1})\,dt\); the fundamental theorem of
calculus in \(L^2\) proves the last assertion.
\end{proof}

\begin{proof}[Proof of the left-density assertion in
Proposition~\ref{prop:ray-lift}]
At a shift \(g\), the left density has the coefficients of \(YY^*\)
weighted, term by term, by the overlap of the two rays based at
\(o,go\), towards \((gh)^+,gh^+\).  There are no such coefficients
unless \(L=r(g)\le4{d_1}+C\): both ray functions are supported within
distance \(2{d_1}+C\) of their basepoints.

We may replace the second endpoint by \((gh)^+\), with an error
whose coefficient \(L^2\) norm is at most
\begin{equation}\label{eq:ray-endpoint-error}
 C e^{C'{d_1}-l}\|Y\|_\infty\|Y\|_2.
\end{equation}
To verify this uniform bound, regard \(h\) in compact polar coordinates
with the additional unary variable \(e^{l-r(h)}\in[0,1]\).
Equation~\eqref{eq:ray-collar} gives smooth agreement of its
current-point direction and ideal endpoint at height zero.
Acting by \(g\), whose length is at most \(4{d_1}+C\), preserves this
vanishing and multiplies each fixed derivative bound by at most
\(e^{C_m {d_1}}\).  One way to see the latter bound directly is to
compose at most \(4{d_1}+C+1\) isometries of length at most one and
differentiate a fixed finite number of times; compact coordinate
transitions have bounded derivatives, and the resulting bounds grow
at most exponentially.  The half-space charts needed at interior
depth at most \(C d_1\) have denominators bounded below by
\(e^{-C_m {d_1}}\), as is also explicit in
\eqref{eq:ray-inversion}.  Formula~\eqref{eq:ray-function}
and the volume on its support have the same fixed-order exponential
bounds.  The overlap difference therefore has a unary smooth
extension with seminorm at most \(C_m e^{C_m {d_1}-l}\), on each of a
fixed finite set of angular charts and the normalized height
interval.  Fourier separation and
Lemma~\ref{lem:ray-unary}, with the other factor uncut,
give \eqref{eq:ray-endpoint-error}.  There is no summation
over the number of shifts in \(\{r\le4{d_1}+C\}\).

For the resulting common-endpoint coefficients, apply
Lemma~\ref{lem:ray-overlap}.  At each annular depth, the
error is a unary symbol of seminorm \(O({d_1}^{-1})\) on a bounded
number of reference boxes around the pole for that \(g\).  Thus its
coefficient \(L^2\) norm is at most
\(C{d_1}^{-1}\|Y\|_\infty\|Y\|_2\), by
Lemma~\ref{lem:ray-unary}.  Telescope the constants
\(a_{d_1}(j,L)\) against single smooth balls.  Their increments have
the same \(O({d_1}^{-1})\) bound and are treated by the same unary
estimate.  The baseline simply multiplies the uncut coefficient by
a bounded scalar.  There are at most \(C(1+{d_1})\) levels because
\(L\le4{d_1}+C\).  Summing these norms and
\eqref{eq:ray-endpoint-error} proves the asserted uniform
bound for the left density.

We record the stronger short-coefficient conclusion needed below.
If \(a=o({d_1})\), then on the shifts \(r(g)\le a\) the baseline
differs from one by \(O(a/{d_1})\), and only \(C(1+a)\) levels occur.
The left-density coefficients restricted to these shifts therefore
differ from those of \(YY^*\) by a quantity bounded in \(L^2\) by
\begin{equation}\label{eq:ray-short-density}
 C\left(\frac{1+a}{d_1}+e^{C'{d_1}-l}\right)
                        \|Y\|_\infty\|Y\|_2.
\end{equation}
This tends to zero when both displayed norms remain bounded, in the
scale regime of Proposition~\ref{prop:ray-lift}.
For cross densities of two bounded long columns the corresponding
statement follows by polarization; the sums and differences of
the columns retain the same long support.  In particular,
\begin{equation}\label{eq:ray-left-pairing}
 \langle\mathbb V_{d_1}Z,\Delta(x)\mathbb V_{d_1}Y\rangle
                               -\langle Z,xY\rangle\longrightarrow0
\end{equation}
uniformly when the operator and Hilbert norms of the long columns
\(Y,Z\), and the Hilbert norm of \(x\), are bounded, with
\(\supp\widehat x\subset\{r\le a\}\).  Explicitly, polarization
and Cauchy--Schwarz bound the absolute value of the difference by
\[
 C\left(\frac{1+a}{d_1}+e^{C'd_1-l}\right)\|x\|_2
   (\|Y\|_\infty+\|Z\|_\infty)(\|Y\|_2+\|Z\|_2).
\]
The Cauchy--Schwarz step pairs the restricted density coefficients
with the Fourier coefficients of \(x\).  In particular it does not
replace \(\|x\|_2\) by an operator bound on an algebra whose trace
mass could vary.

For infinite Fourier support, the bounded approximants with slacked
radial cutoffs used after \eqref{eq:input-7} apply here as well.
The density bounds yield weakly convergent \(L^2\) subsequences,
whose trace pairings identify their limits with the vector densities.
This proves (i) in full generality.
\end{proof}

\begin{proof}[Proof of the pairing and intertwining assertions in
Proposition~\ref{prop:ray-lift}]
First prove (iii).  The kernel \(K_{d_1}(v,w)\) satisfies
\(K_{d_1}(v,v)=1\).  For each fixed derivative order \(m\), its
derivatives on fixed angular charts are bounded by
\(C_m e^{C_m {d_1}}\): use smooth compact frames, the depth restriction
\(d(o,p)\le2{d_1}+C\), and the same fixed-order coordinate and volume
bounds as in \eqref{eq:ray-endpoint-error}.  On a pair of
adjacent boxes of radius \(e^{-d_2}\), Taylor's formula on the
diagonal and the weighted coordinate scaling imply
\begin{equation}\label{eq:ray-kernel-near-one}
 \|K_{d_1}-1\|_{C^m_{\rm normalized}}
                  \le C_m e^{C_m {d_1}-d_2}.
\end{equation}
For derivatives of positive order, at least one normalized angular
derivative supplies a factor \(e^{-d_2}\); the zero-order estimate
uses the equality on the diagonal.  Central derivatives supply
the stronger factor \(e^{-2d_2}\).

Choose \(d_2\) with \({d_1}\ll d_2\ll l'\).  The assumed support
condition makes every contributing pair adjacent at scale
\(e^{-d_2}\).  Insert a smooth partition at that scale and separate
\(K_{d_1}-1\) on each adjacent pair.  At one frequency, Cauchy--Schwarz
for the columns of cut vectors gives
\[
 \left|\sum_{C,D'}t_{C,D'}\langle Z_C,Y_{D'}x\rangle\right|
 \le\|t\|\,\|x\|_\infty
       \left(\sum_C\|Z_C\|_2^2\right)^{1/2}
       \left(\sum_{D'}\|Y_{D'}\|_2^2\right)^{1/2}.
\]
The matrix norm is bounded by its degree times its largest
coefficient, which tends to zero by
\eqref{eq:ray-kernel-near-one}.  Overlap bounds control
the two sums by \(C\|Z\|_2^2,C\|Y\|_2^2\).
Choose the finite derivative order in advance sufficiently large
for a summable Fourier expansion.  Since \(d_2/{d_1}\to\infty\), the
sum of the frequency bounds still tends to zero.  This proves
\eqref{eq:ray-close-pairing}.  The proof only needs
\(Y,Z\in L^2\); Fourier truncation of those two vectors and
continuity extend it with \(x\) remaining bounded.

For right multiplication by a short element we verify the support
condition.  If \(r(h),r(hk)\ge l-a\) and \(r(k)\le a\), then
\(d(ho,hko)\le a\).  In a compact finite boundary chart the
heights of these two points are comparable to their radial
exponentials.  Their horizontal separation, measured by a
homogeneous gauge, either is of that small height order or, by
\eqref{eq:ray-distance}, satisfies
\[
 d(ho,hko)\ge r(h)+r(hk)
       +2\log d_{\partial X}(h^+,(hk)^+)-C.
\]
The collar errors are absorbed in the first alternative.  A finite
chart cover, or a compact rotation placing the pair in one such
chart, gives
\(d_{\partial X}(h^+,(hk)^+)\le Ce^{-c(l-2a)}\)
with fixed \(c>0\).  Since \(a/{d_1}\to0\) and \(l/{d_1}\to\infty\),
assertion (iii) applies.  It also applies with \(xx^*\), supported
in \(\{r\le2a\}\), and with the long column \(Yx\).
Expanding the square of
\(\mathbb V_{d_1}(Yx)-(\mathbb V_{d_1}Y)x\), its three terms tend respectively
to \(\|Yx\|_2^2,\|Yx\|_2^2,-2\|Yx\|_2^2\).
This proves the right assertion in (ii).

For the left assertion use
\eqref{eq:ray-left-pairing} first with \(x^*x\), then with
the cross pair \(Z=xY\) and \(x\).  The support of \(x^*x\) lies
in \(\{r\le2a\}\), and \(xY\) has support in
\(\{r\ge l-a\}\).  The same expansion of the squared norm now
proves \(\|\mathbb V_{d_1}(xY)-\Delta(x)\mathbb V_{d_1}Y\|_2\to0\).

All proofs used trace identities, Plancherel, and constants uniform
in the coefficient algebra.  They therefore apply to matrix
algebras and, after integration, to measurable fields.  To make
any support-preserving approximation needed here, first use bounded
Fourier-polynomial approximation and then smooth radial cutoffs
with fixed slack.  Proposition~\ref{prop:ray-radial} gives
uniform operator bounds; doubling a short support and halving a
long support preserves all scale separations in the statement.
\end{proof}

\begin{corollary}[Stability under Hilbert-small changes]
\label{cor:ray-small-change}
If \(Y\) is a long column as above with
\(\|Y\|_\infty\) and \(\|Y\|_2\) uniformly bounded and
\(x-x'\to0\) in \(L^2(D)\), with \(x,x'\) uniformly bounded, then
\[
 \|\Delta(x-x')\mathbb V_{d_1}Y\|_2\longrightarrow0,
 \qquad \|(\mathbb V_{d_1}Y)(x-x')\|_2\longrightarrow0.
\]
The analogous statement holds for bounded elements in the amplified
left algebra when their difference tends to zero in its semifinite
\(L^2\) norm.
\end{corollary}

\begin{proof}
Pair the appropriate density from Proposition~\ref{prop:ray-lift}
with \((x-x')^*(x-x')\), or with the reversed product on the right.
Tracial Cauchy--Schwarz and
\[
 \|(x-x')^*(x-x')\|_2
             \le\|x-x'\|_\infty\|x-x'\|_2
\]
give, with \(z=x-x'\), the homogeneous bound
\begin{equation}\label{eq:ray-small-change-bound}
 \max\{\|\Delta(z)\mathbb V_{d_1}Y\|_2^2,
              \|(\mathbb V_{d_1}Y)z\|_2^2\}
 \le C\|Y\|_\infty\|Y\|_2\|z\|_\infty\|z\|_2.
\end{equation}
In the amplified case use the density
\(\mathbb V_{d_1}Y(\mathbb V_{d_1}Y)^*\) and the same inequality
with the semifinite Hilbert norm of \(z\).
For measurable fields, integrate this inequality.  Uniform
operator bounds and Cauchy--Schwarz on the base give the product of
the integrated Hilbert norms of \(Y\) and \(z\).  Thus integrated
Hilbert-small changes suffice; no pointwise error bound is required.
\end{proof}

\subsection{Admissible representations and spherical laws}

Fix an isolated measured factor \(T\) of \(G\), its symmetric space
\(X_T\), and a base point \(o_T\).  Write
\[
 r_T(h)=d(o_T,h_T o_T),\qquad h\in\Gamma,
\]
with distance normalized by the indivisible simple root.  The
projection \(h\mapsto h_T\) need not have discrete image.  In this
Section \(D_f\) denotes the Fourier multiplier with symbol \(f\).
All finite factor traces are normalized.

\begin{definition}\label{def:ray-admissible}
An admissible representation for \(T\) consists of a finite factor
\[
 D=B\bar\otimes L_\eta(\Gamma)
\]
and a normalized projective representation
\(\Theta:\Gamma\to\mathcal U(D)\) such that
\[
 D\cap\Theta(\Gamma)'=\mathbb C1
\]
and the representation
\begin{equation}
 \gamma\longmapsto\Ad\Delta_T(\Theta(\gamma))
 \quad\hbox{on}\quad
 L^2\bigl(B(L^2(X_T))\bar\otimes D,\Tr\otimes\tau\bigr)
 \label{eq:ray-amplified-action}
\end{equation}
has no invariant vector.  Here
\(\Delta_T(u_h)=\lambda_{X_T}(h_T)\otimes u_h\), with the identity
action on \(B\).  The scalar multiplier of \(\Theta\) is not required
to equal \(\eta\).
\end{definition}

Keep the section \(d_0:Y=G/\Gamma\to G\) and cocycle \(c\) from
Section~\ref{sec:setup}.  For \(a\in G\), put
\[
 X^\Theta_a(y)=\Theta(c(a,y)).
\]
If \(S\) is an isolated source factor and \(a_t\in S\) is a positive
split translation of length \(t\), the spherical law at \(t\), on
scale \(R>0\), is the probability measure
\begin{equation}
 P_{\Theta,R}(t)(E)=
 \int_Y\bigl\|D_{1_{\{r_T/R\in E\}}}X^\Theta_{a_t}(y)\bigr\|_2^2\,dy
 \quad (E\subset[0,\infty]\text{ Borel}).
 \label{eq:ray-spherical-law}
\end{equation}
The value \(+\infty\) is used only in weak limits.  Addition on
\([0,\infty]\) is the usual extended addition.

\begin{lemma}[Uniform short approximation]\label{lem:ray-compactshort}
Let \(R_n\to\infty\), and let \(\Theta_n\) be projective
representations in finite tensors as above.  Suppose
\begin{equation}
 \|D_{1_{\{r_T>\varepsilon R_n\}}}\Theta_n(\gamma)\|_2\longrightarrow0
 \quad(\gamma\in\Gamma,\ \varepsilon>0).
 \label{eq:ray-short-fixed}
\end{equation}
There are \(a_n=o(R_n)\) and uniformly operator-bounded fields
\(Z_{n,s}\), with Fourier support in \(\{r_T\le a_n\}\), such that
\begin{equation}
 \sup_{s\in C}\|X^{\Theta_n}_s-Z_{n,s}\|_{L^2(Y;L^2(D_n))}
 \longrightarrow0
 \label{eq:ray-short-compact}
\end{equation}
for every compact \(C\subset G\).  Consequently a compact family of
source translations before or after a segment does not change any
limiting rescaled length law of that segment.
\end{lemma}

\begin{proof}
Choose a compact \(K\subset G\) for which
\(|\{y:d_0(y)\notin K\}|<\varepsilon\).
Uniformly for \(s\in C\), outside a set of measure at most
\(2\varepsilon\), both \(d_0(y)\) and \(d_0(sy)\) belong to \(K\).
There
\[
 c(s,y)\in \Gamma\cap K^{-1}CK,
\]
a finite set.  Apply \eqref{eq:ray-short-fixed} to this finite
set.  Diagonalization over compact exhaustions of \(G\), followed
by Proposition~\ref{prop:ray-radial}, gives
\eqref{eq:ray-short-compact}.  More explicitly, use a fixed
smooth function equal to one on \([0,1]\) and zero on \([2,\infty)\),
at a cutoff \(b_n=o(R_n)\) tending sufficiently slowly to infinity
relative to the required finite lists, and put \(a_n=2b_n\).
The completely bounded multiplier norm gives the uniform
operator bound.  The normalized factor traces also give
\(\|X_s^{\Theta_n}(y)\|_2=1\); Plancherel bounds the integrated
Hilbert norms of the short approximants by the supremum norm of
the fixed cutoff.  Thus both kinds of bounds required in
Proposition~\ref{prop:ray-lift} are uniform in these applications.

Here is the spectral assertion needed in the last sentence.
Multiplication by an element supported in \(\{r_T\le a_n\}\) sends
the Fourier support of a radial interval into its
\(a_n\)-neighborhood, on either side, because
\begin{equation}
 |r_T(gh)-r_T(h)|\le r_T(g),\qquad
 |r_T(hg)-r_T(h)|\le r_T(g).
 \label{eq:ray-triangle}
\end{equation}
Take two closed rescaled intervals at positive distance.  A short
multiplier cannot connect their radial cuts for large \(n\).
For smooth radial cuts of a uniformly bounded column the errors
from replacing an actual transporter by its short approximation
are Hilbert-small: one factor is Hilbert-small and the other is
uniformly operator-bounded.  The radial multiplier bound supplies
the latter bound for every cut used here.  Finite spectral
partitions therefore show that multiplication by the actual
unitary transporter intertwines the limiting rescaled length
measurement.  Approximation of continuous functions on
\([0,\infty]\) removes the partitions.  Reindexing \(Y\) preserves
its probability measure, and the scalar phases in a product of
transporters cancel from all these Gram forms.
\end{proof}

\begin{lemma}[Averaging with a uniform gap]\label{lem:ray-kazhdan}
There are a compactly supported, absolutely continuous, symmetric
probability \(m_0\) on \(G\) and \(0<c_0<1\) such that its
convolution operator is positive in every unitary representation
and has norm at most \(c_0\) on the orthogonal complement of the
invariant vectors.
\end{lemma}

\begin{proof}
Let \((C,\varepsilon)\) be a compact Kazhdan pair for the whole
locally compact product \(G\); property~\((T)\) provides such a
pair \cite{BHV2008}.  Choose a
relatively compact neighborhood \(V\) of the identity with positive
Haar measure, and a symmetric compactly supported probability
density \(\nu\) bounded below on \(V\cup CV\) and their inverses.
For a unit vector \(\xi\), writing \(g=(gv)v^{-1}\) and averaging
over \(v\in V\) bounds
\[
 \sup_{g\in C}\|\pi(g)\xi-\xi\|^2
 \le C_1\int_G\|\pi(h)\xi-\xi\|^2\,d\nu(h).
\]
Indeed the two terms in the triangle inequality are integrated
over \(CV\) and \(V^{-1}\), where the density has a uniform
positive lower bound.  The Kazhdan inequality thus gives a
uniform upper bound \(1-\delta\) for the self-adjoint averaging
operator \(\pi(\nu)\) on the complement of the invariant vectors.
The lazy average \((1+\pi(\nu))/2\) is positive and has norm at
most \(1-\delta/2\) there.  Convolve this probability on both sides
with a symmetric compactly supported probability density \(\chi\).
Its operator is
\[
 \pi(\chi)\frac{1+\pi(\nu)}2\pi(\chi),
\]
which is positive and has the same uniform upper bound.
The resulting probability \(m_0\) is absolutely continuous.
\end{proof}

\begin{proposition}[Addition in distribution]\label{prop:ray-addition}
Let \(R_n\to\infty\), let \(\Theta_n\) be admissible for \(T\), and suppose
\eqref{eq:ray-short-fixed} holds.  Fix an isolated source
factor and a positive split axis in it.  For any \(t_n,u_n\ge0\),
if the individual spherical laws converge weakly to \(P_t,P_u\),
then
\begin{equation}\label{eq:input-8}
 P_{\Theta_n,R_n}(t_n+u_n)\longrightarrow P_t*P_u .
\end{equation}
The same statement permits pre- and post-translations in fixed
compact subsets of \(G\).
\end{proposition}

\begin{proof}
\par\smallskip\noindent\textit{The inserted path and the order of limits.}\quad
First pass to a subsequence on which the sum laws converge.
Insert a fixed translation \(s\in G\) between the two positive
axial segments.  Away from the Bruhat wall corresponding to
backtracking, Cartan multiplication expresses
\(a_{t_n}s a_{u_n}\) as \(a_{t_n+u_n}\) with pre- and
post-translations in a compact set depending on \(s\).
This includes a segment whose length stays bounded.  To verify
the assertion when both lengths diverge, write \(s\) in the
open opposite-radical--Levi--radical cell for the axis: the two
radical factors contract at the respective ends, and the Levi
factor has a fixed split component.  The omitted wall is a
proper lower-dimensional Bruhat cell and has Haar measure zero.
The components of \(s\) outside the chosen source factor are
fixed compact displacements.  The exceptional wall in that factor
has Haar measure zero, also after taking its product with the
complementary factors.  Lemma~\ref{lem:ray-compactshort} therefore identifies the uncut law
of the inserted product with the limiting sum law, for almost
every \(s\).

\par\smallskip\noindent\textit{The Gram computation.}\quad
Write \(A_n(y)=X^{\Theta_n}_{a_{t_n}}(y)\) and
\(B_n(y)=X^{\Theta_n}_{a_{u_n}}(a_{u_n}^{-1}y)\).
For smooth radial functions \(f,g\), constant near the two
ends of the compactified interval, put
\begin{equation}
 A_{n,f}=D_{f(r_T/R_n)}A_n,\qquad
 B_{n,g}=D_{g(r_T/R_n)}B_n .
 \label{eq:ray-cut-columns}
\end{equation}
Consider the columns
\begin{equation}
 A_{n,f}(sy)X^{\Theta_n}_s(y)B_{n,g}(y).
 \label{eq:ray-inserted-columns}
\end{equation}
They are uniformly operator-bounded.  Their mixed Gram forms
are pairings between
\(A_{n,f'}^*A_{n,f}\) and the translate of
\(B_{n,g}B_{n,g'}^*\) under the induced trace action
\[
 F(y)\longmapsto
 \Ad\Theta_n(c(s,s^{-1}y))F(s^{-1}y).
\]
Its invariant space consists of constant scalar fields:
an invariant equivariant field lifts to a constant field on
\(G\), whose value commutes with \(\Theta_n(\Gamma)\).
More explicitly, the lifted field is left invariant, hence
essentially constant by Haar invariance; its residual right
\(\Gamma\)-equivariance imposes precisely that commutation
relation.  Conversely a scalar constant gives an invariant field.
The matrix products appearing in the Gram forms have uniformly
bounded integrated Hilbert norms, since each radial piece has
both a uniform operator bound and a uniform Hilbert bound.
Average \(s\) with \(m_0^{*k}\).  Lemma~\ref{lem:ray-kazhdan}
shows that the difference from the product of the scalar Gram
forms is \(O(c_0^k)\), uniformly in \(n\).

Take first a Hilbert limit in \(n\), for each fixed \(k\), and
then a Hilbert limit in \(k\).  Equivalently, form the two
successive Hilbert ultraproducts of the spaces with measures
\(dy\,dm_0^{*k}(s)\).  Polarization of the preceding calculation
gives an isometry
\begin{equation}
 J:L^2(P_t)\otimes L^2(P_u)\longrightarrow\mathcal K,\qquad
 f\otimes g\longmapsto
 \bigl(A_{n,f}(sy)X^{\Theta_n}_s(y)B_{n,g}(y)\bigr)^{\lim}.
 \label{eq:ray-convolution-link}
\end{equation}
The Fourier length in the output has its own spectral
measurement \(E\) on \(\mathcal K\).
For each fixed \(k\), the measure \(m_0^{*k}\) has compact support.
All short-approximation errors may therefore be sent to zero in
\(n\) before increasing \(k\).  Dominated convergence identifies
the uncut output law for that fixed average with the subsequential
sum law.  No uniform approximation over the growing supports of
all convolution powers is asserted or needed.

\par\smallskip\noindent\textit{Positive finite input windows.}\quad
We prove that its restriction to the range of \(J\) is addition.
An output interval strictly above the sum of two input
intervals is excluded by \eqref{eq:ray-triangle} and short
approximation of \(X^{\Theta_n}_s\).  It remains to exclude a
strict deficit.  Localize the two input lengths in smooth
windows bounded below by \(\varepsilon'R_n\), and localize the
output in a window shorter than their sum by
\(\varepsilon R_n\), with \(\varepsilon,\varepsilon'>0\).
All window boundaries may be enlarged slightly.  Write
\(A=A_{n,f}\), \(B=B_{n,g}\).  In testing the output cut,
trace cyclicity gives pairings of \(A(sy)^*\) against
\[
 X^{\Theta_n}_s(y)B(y)C(y,s),
\]
where \(C\) is uniformly operator-bounded and has the Fourier
support of the inverse shortened-output window.  In particular
one may take \(C\) to be the adjoint of the smooth squared output
cut of the product itself.

Choose, by Lemma~\ref{lem:ray-compactshort}, scales
\[
 a_n\ll d_n\ll R_n .
\]
On each compact set of \(s\)'s replace \(X^{\Theta_n}_s\) by
its short approximation.  Lemma~\ref{lem:ray-defect} implies that every contributing
pair of long indices has directions at distance at most
\(C e^{-c\varepsilon R_n}\).  Proposition~\ref{prop:ray-lift}
therefore permits insertion of the two ray lifts and then
intertwines the short left multiplication through the lift.
The result, with an error tending to zero, is
\begin{equation}
 \tau\bigl(\zeta_n(sy)^*
       \Delta_T(X^{\Theta_n}_s(y))\xi_n(y)C(y,s)\bigr),
 \quad
 \xi_n=\mathbb V_{d_n}B,\quad
 \zeta_n=\mathbb V_{d_n}(A^*).
 \label{eq:ray-shortening-test}
\end{equation}
The errors are uniform for \(s\) in a fixed compact set.
Replacing the short approximation by the actual transporter
on the lifted vector is legitimate because the left density
has uniformly bounded tracial \(L^2\)-norm: for a bounded
Hilbert-small change \(z\),
\begin{equation}
 \|\Delta_T(z)\xi_n\|^2
 \le \|z\|_\infty\|z\|_2\|\rho^{\,L}_{\xi_n}\|_2 .
 \label{eq:ray-density-error}
\end{equation}
The integrated version follows from Cauchy--Schwarz on \(Y\).

The exact trace inequality for \eqref{eq:ray-shortening-test}
is
\begin{align}
 &\left|\tau\bigl(\zeta_n(sy)^*
       \Delta_T(X^{\Theta_n}_s(y))\xi_n(y)C(y,s)\bigr)\right|^2
 \nonumber\\
 &\quad\le
 \|C(y,s)\|_2^2
 (\Tr\otimes\tau)\left[
 \zeta_n(sy)\zeta_n(sy)^*
 \Ad\Delta_T(X^{\Theta_n}_s(y))
             \bigl(\xi_n(y)\xi_n(y)^*\bigr)\right].
 \label{eq:ray-density-pairing}
\end{align}
Both density fields on the right have uniformly bounded
trace-Hilbert norm, by Proposition~\ref{prop:ray-lift};
they are independent of \(s\) apart from the displayed transport.
The induced amplified action has no invariant vector, by
admissibility and the same equivariant-field lifting argument.
Its \(m_0^{*k}\)-average is therefore bounded by \(C c_0^k\).
Integrate \eqref{eq:ray-density-pairing}, first let \(n\)
tend to infinity, and then let \(k\) tend to infinity.
The shortened-output cut vanishes.

\par\smallskip\noindent\textit{Zero and infinite input values.}\quad
Suppose first that one input value is zero and the other is a
finite value \(v\).  Restrict the first input to \([0,\delta]\)
and the second to \([v-\delta,v+\delta]\cap[0,\infty)\), on the
rescaled length axis.  The short intermediate transporter changes
length by \(o(1)\) on this scale.  The forward and reverse triangle
inequalities therefore restrict the output to
\([v-2\delta-o(1),v+2\delta+o(1)]\cap[0,\infty)\).
Every closed output window separated from \(v\) is excluded by
choosing \(\delta\) sufficiently small.  This includes \(v=0\),
where the upper triangle bound suffices.  The cut replacements have
the same Hilbert-small errors as before.

For an input value \(+\infty\), it suffices to exclude output
in a fixed finite interval \([0,M]\).  If the other input is
finite, choose a bounded neighborhood \([0,B]\) containing it
and restrict the infinite input to \([N,\infty]\), with
\(N>M+B+1\).  The reverse triangle inequality excludes the
output window, after the short intermediate replacement.  If
both input values are infinite, restrict both to
\([N,\infty]\), where \(2N>M+1\).  The connecting output index
then has a deficit at least \((2N-M-o(1))R_n\) from the sum of
the two long input lengths.  The lifted shortening argument just
proved applies: it used a positive lower length bound and a
uniform deficit, but no upper bound on either long input.
Use smooth tail cutoffs constant near infinity, whose multiplier
norms are bounded by Proposition~\ref{prop:ray-radial}.  Hold
\(N,M\) fixed through the limits in \(n\) and then \(k\);
afterward exhaust the finite output intervals and increase the
finite lower input cutoffs.  This excludes every finite output
value at an infinite sum.

\par\smallskip\noindent\textit{Identification of the output measurement.}\quad
We have proved that a product of supported smooth input functions
has no output in a closed window disjoint from the sums of its
input supports.  Characteristic functions of interval partitions
whose boundaries have zero input measure are obtained by Hilbert
approximation with such smooth functions.  The isometry \(J\)
preserves this approximation.

Here is an exact spectral argument that does not assume in advance
that \(E\) preserves the range of \(J\).  Write
\(a(s,t)=s+t\), a continuous map on the compact square, and let
\(P_{\rm in}\) be multiplication by Borel subsets of that square
on \(L^2(P_t)\otimes L^2(P_u)\).  Fix a closed output set \(F\).
Product test functions supported on rectangles whose closures lie
in the open set \(a^{-1}(F)^c\) have dense span in
\(\ran P_{\rm in}(a^{-1}(F)^c)\), by regularity of the product
measure and Hilbert approximation with smooth interval cutoffs.
For each such rectangle the compact sum support misses \(F\).
A finite cover of \(F\) by the excluded output windows and the
rectangle support conclusion therefore give
\[
 E(F)J P_{\rm in}(a^{-1}(F)^c)=0,
 \qquad
 E(F)J=E(F)J P_{\rm in}(a^{-1}(F)).
\]
Choose closed sets \(K_j\uparrow F^c\); for example use positive
distance at least \(1/j\) from \(F\) in a compatible compact
metric.  Apply the same conclusion to every \(K_j\).
Since \(a^{-1}(K_j)\) is disjoint from \(a^{-1}(F)\), it gives
\[
 E(K_j)J P_{\rm in}(a^{-1}(F))=0.
\]
Strong convergence of these spectral projections yields
\(E(F^c)J P_{\rm in}(a^{-1}(F))=0\).  Combining the two identities
proves
\[
                    E(F)J=J P_{\rm in}(a^{-1}(F)).
\]
The monotone-class theorem extends this equality to all Borel
sets, and spectral integration then gives
\[
 E(\varphi)J=J\,M_{\varphi(s+t)}
 \qquad(\varphi\in C([0,\infty])).
\]
On \(J(1\otimes1)\), the output law is the limiting sum law
identified in the first paragraph.  This proves
\eqref{eq:input-8}.  Compact pre- and post-translations
are covered by Lemma~\ref{lem:ray-compactshort}.  Every convergent
subsequence of the sum laws has the same limit, so compactness of
the space of probability measures on \([0,\infty]\) proves
convergence of the original sequence.
\end{proof}
